\section{Introduction}
\label{sec:introduction}

\paragraph{Motivation.}
A central trend in large language model deployment is the explicit allocation
of test-time compute to deliberation. Modern reasoning systems such as
OpenAI's o1 line \citep{openai2024o1} and the public scaling analyses of
\citet{snell2024scaling} demonstrate that, for verifiable problems,
additional inference-time computation can yield accuracy gains comparable to
those produced by orders-of-magnitude increases in pretraining scale. At the
same time, small open-weight models, on the order of one to seven billion
parameters, lag far behind frontier reasoning models on mathematical and
programmatic benchmarks such as GSM8K \citep{cobbe2021gsm8k} and MATH
\citep{hendrycks2021math}, even when given generous decoding budgets. The
gap is widest precisely where it should be most closable: on problems whose
solutions can be checked by a non-LLM verifier such as a symbolic algebra
system, a unit-test harness, or a proof assistant. Process-level
verifier supervision has been shown to outperform outcome-only supervision
on MATH \citep{lightman2024verify}, evidence that verifier signals carry
usable learning content at fine granularity. Whenever a cheap and
trustworthy verifier exists, the loop from candidate solution to scalar
reward is also short and dense, and, because the reward is a deterministic
check rather than a learned critic, it sidesteps the reward-hacking surface
that a learned proxy opens under optimization
\citep{skalse2022reward, gao2023scaling}. This is the classical setting in
which reinforcement learning excels.

We therefore frame the problem squarely as reinforcement learning. The agent
interacts with an environment whose dynamics are induced by a base language
model and whose terminal reward is supplied by an external verifier. Learning
proceeds by approximate $Q$-learning over sampled transitions, in the
direct lineage of \citet{sutton2018reinforcement} and the deep $Q$-learning
techniques of \citet{mnih2015human}. The natural question is then whether a
small base model, paired with a learned action-value function trained against
verifier-anchored Bellman targets, can close a meaningful fraction of the gap
to much larger reasoning models on a fixed compute axis. This paper is the
methodological prerequisite to answering that question.

\paragraph{Position.}
This is a methodology paper that presents a plan, not results. We make the
framework precise enough to be evaluated as a contribution in its own right, and
we pre-register the experimental design so that the eventual benchmark study
will not be confused with post-hoc rationalization. We also report a preliminary
feasibility probe (Section~\ref{sec:pilot}) that exercised the full pipeline
end-to-end; the probe predates the current action space, so we report it only as
evidence that every component is implementable and as the source of several
methodological precautions the plan inherits, not as a benchmark result and not
as a number the framework's claims rest on. Concretely, we present a
reinforcement-learning framework specialized for verifiable LLM-induced context
Markov decision processes, an architecture that uses a shared frozen language
model as both encoder and stochastic transition model, a reward structure that
quarantines learned aesthetic signals away from the Bellman target, an explicit
role for tree search as a policy improvement operator rather than as the central
learning algorithm, and a pre-registered evaluation plan with falsifiable
hypotheses against standardized benchmarks.

The reframing relative to prior MCTS-centric work is deliberate. In the
AlphaZero tradition \citep{silver2017alphagozero, silver2018alphazero} and its
expert-iteration ancestor \citep{anthony2017thinking}, tree search is the
narrative center and the learned policy is its student. In our setting the
center is the agent and its action-value function $\Qpsi(s, a)$. Tree search
is one of several mechanisms for producing better samples and stronger
training targets. This shift in emphasis has practical consequences: it makes
non-search baselines such as $\epsilon$-greedy and softmax exploration first-class
citizens, it clarifies what the verifier is for, and it makes the test-time
search knob optional rather than load-bearing.

\paragraph{Contributions.}
We make six contributions:
\begin{enumerate}[leftmargin=*]
  \item A reinforcement-learning framework that casts LLM deliberation on
        verifiable tasks as an approximate $Q$-learning problem over an
        induced context MDP, with a structured state schema
        (Section~\ref{sec:method-state}) that makes the Markov bet explicit
        and a canonical \textsc{Verify} action that is part of
        the agent's action space rather than a post-hoc parser pass.
  \item An architecture that uses a single frozen base language model as both
        encoder and stochastic transition model, with a small adapter and a
        multi-head parametric action-value network $\Qpsi(s, a)$ over a fixed
        verb set (Section~\ref{sec:method-architecture}). The choice of a
        multi-head $Q$ over a fixed action set, rather than a scoring head over
        $(s, a)$ embeddings, is justified by the available action cardinality
        and by the gradient-routing properties of the architecture.
  \item A reward-design principle that separates verifier-anchored Bellman
        targets from learned-aesthetic signals
        (Section~\ref{sec:method-reward}). Verifier outcomes form the external
        reward and the Bellman backup; learned signals such as
        reasoning-quality classifiers or deception detectors may enter only as
        search priors, never as reward, reducing the Goodhart attack surface.
  \item Monte Carlo Tree Search as a sample-efficient policy improvement
        operator (Section~\ref{sec:method-mcts}), producing stronger training
        targets $\Qmcts(s, a)$ and a visit-count policy $\pimcts$ that are
        distilled into the parametric $Q$ and policy heads. The same operator
        serves as a test-time-compute knob at inference: a fast greedy mode
        evaluates $\Qpsi$ once per state, while a search mode invokes MCTS
        when latency permits. Search is therefore an optional accelerator,
        not the central learning algorithm.
  \item A pre-registered evaluation plan
        (Section~\ref{sec:experimental-plan}) that defines the baseline grid,
        the compute axis, and the headline compute-versus-accuracy frontier
        figure before any training begins, with falsifiable hypotheses and an
        explicit threat model for compute leaks in test-time-compute claims.
  \item A preliminary feasibility probe (Section~\ref{sec:pilot}) that exercised
        every component of the framework end-to-end and established that the
        expert-iteration loop is implementable, while surfacing the
        methodological precautions the pre-registered study inherits as
        monitoring requirements: train/held-out structural discipline, per-class
        composition tracking, a pairwise-recommended-verb target-stability
        diagnostic with the budget set by the slowest-converging family, and
        honest per-decision compute accounting with an observable, resumable run
        harness. The probe predates the current action space, so it contributes
        feasibility evidence and methodology, not benchmark numbers.
\end{enumerate}

\paragraph{Paper structure.}
Section~\ref{sec:background} develops the reinforcement-learning background:
Markov decision processes, the Bellman optimality equation, tabular and
approximate $Q$-learning with target networks, exploration policies, and MCTS
as a policy improvement operator in the AlphaZero tradition, ending with the
LLM-induced context MDP framing. Section~\ref{sec:method} presents the
framework in detail: the state schema, the action space including the verifier
action, the shared-encoder architecture, the role of MCTS in producing
training targets and at inference, the reward structure, the training loop
combining replay-based approximate $Q$-learning with expert-iteration
distillation, and the two inference modes.
Section~\ref{sec:experimental-plan} specifies the experimental design and the
hypotheses to be tested. Section~\ref{sec:pilot} reports the preliminary
feasibility probe: what it exercised end-to-end, the methodological lessons that
survive the action-space redesign, and the design implications carried forward
into the pre-registered study.
Section~\ref{sec:related-work} positions the contribution against prior work
in classical RL, verifier-supervised reasoning, test-time-compute scaling,
MCTS-augmented language models, and small-model self-improvement.
Section~\ref{sec:discussion} discusses limitations, extensions, and threats
to validity. Section~\ref{sec:conclusion} concludes.
